\section{An exact smooth correction}
\label{sec:exact-correction}

The small increment of Proposition~\ref{prop:small-increment} leaves a
metric error.  We now remove that error while keeping the immersion
close in $C^2$.  Smoothing is used only to construct each increment;
the current map itself is never replaced by its smoothing.  This
preserves its higher derivatives until a sufficiently late stage can
control them.

Smoothing in an iteration with derivative loss originates in Nash's
smooth embedding argument \cite[Part~A]{Nash1956} and was developed by
Moser \cite[Chapter~I, \S\S1 and~5]{Moser1966}; Hamilton gives a systematic
inverse-function formulation \cite[Part~III]{Hamilton1982}.
Here we establish the correction estimates and the passage through the
derivative orders directly, without assuming a smooth family of inverses
of the full linearized metric map.

\begin{proposition}[Exact correction]\label{prop:exact-correction}
Let $M$ be a closed smooth surface, and let $h_*$ be a smooth Riemannian
metric on $M$.  Fix background norms.  There are positive integers
$R,N$ with the following property.  Suppose that $f_z:M\to\R^4$ are
smooth immersions for sufficiently small $z>0$, and that
\begin{align}
 \|f_z\|_{C^m}&=O(z^{-m}) &&(0\le m\le R+2),\label{eq:exact-input-jets}\\
 \|\gamma(f_z)-h_*\|_{C^R}&=O(z^N),\label{eq:exact-input-metric}\\
 \min_{p\in M}\max_{|u|=1}|B_{f_z}(u,u)|&\ge c>0.
 \label{eq:exact-input-second-form}
\end{align}
Then, for all sufficiently small $z$, there is a smooth immersion
$\widetilde f_z:M\to\R^4$ such that
\[
 \gamma(\widetilde f_z)=h_*,\qquad
 \|\widetilde f_z-f_z\|_{C^2}=o(1).
\]
The integers $R,N$ do not depend on the constants in
\eqref{eq:exact-input-jets}--\eqref{eq:exact-input-second-form}.
The threshold on $z$ may depend on them.  No regularity of the family
with respect to $z$ is required or asserted.
\end{proposition}

The proof has two parts.  First we obtain a recurrence valid in every
fixed derivative order, from a bound in only one finite order.  We
then let that controlled order increase, waiting a finite number of
steps whenever necessary.  The distinction between these two orders
is what gives a smooth limit from the finite assumptions above.

\subsection{Smoothing with a finite input bound}

We use the weighted norms introduced in the preceding section.
The following elementary smoothing facts also apply to tensor fields.

\begin{lemma}\label{lem:finite-order-smoothing}
For each positive integer $r$, there are linear smoothing operators
$S_s$, $0<s\le s_0$, on smooth sections of any fixed smooth vector
bundle over $M$, with the following properties.  For fixed
nonnegative integers $q,m$, and $0\le j\le r$,
\begin{align}
 \|S_s u\|_{C^{q+m}}&\le C_{q,m,r}s^{-m}\|u\|_{C^q},
 \label{eq:smooth-derivative}\\
 \|u-S_su\|_{C^{q+j}}&\le C_{q,j,r}s^{r-j}\|u\|_{C^{q+r}},
 \label{eq:smooth-approximation}\\
 \|S_su\|_{C^m}&\le C_{m,r}\|u\|_{C^m}.
 \label{eq:smooth-boundedness}
\end{align}
The constants use only the fixed atlas, bundle, and displayed orders.
\end{lemma}

\begin{proof}
Choose a compactly supported smooth kernel $K$ on $\R^2$ with integral
one.  A finite linear combination of its distinct dilates can be
chosen to have integral one and all moments of degrees $1,\ldots,r$
equal to zero.  Indeed, if the dilation factors are distinct positive
numbers $b_0,\ldots,b_r$, choose their coefficients $c_i$ by
\[
 \sum_{i=0}^r c_i b_i^d=
 \begin{cases}1,&d=0,\\0,&1\le d\le r.\end{cases}
\]
The coefficient matrix is a Vandermonde matrix.  Every moment of
multi-degree $\alpha$ scales by $b_i^{|\alpha|}$, so these equations
cancel all the required moments.  The resulting kernel need not be
positive.

Take a finite coordinate and bundle atlas with a subordinate smooth
partition of unity.  Extend each localized component by zero in its
coordinate chart, convolve at scale $s$, and multiply by a second
cutoff equal to one near the original support.  For sufficiently
small fixed $s_0$, all these operations occur within the charts.
Sum the resulting sections.  At scale zero this sum is exactly $u$.
Taylor's formula and the moment conditions give
\eqref{eq:smooth-approximation}.  Moving derivatives between the
localized component and the kernel gives
\eqref{eq:smooth-derivative}; moving all derivatives onto the component
gives \eqref{eq:smooth-boundedness}.  Derivatives of the fixed cutoffs
and bundle frames introduce only lower derivatives with fixed
coefficients.  This proves all three bounds.
\end{proof}

We record how these facts will be used.  Let $G:M\to\R^4$ be smooth,
let $H$ be a smooth symmetric tensor, and define
\begin{equation}\label{eq:exact-weighted-input}
 \mathcal S_m(t;G,H)=
 \sum_{j=0}^m t^j\bigl(\|G\|_{C^{j+2}}+\|H\|_{C^j}\bigr).
\end{equation}
Suppose $\mathcal S_r(t;G,H)\le B$, and let
$0<\tau\le s\le t\le1$.  Set $G_s=S_sG$ and $H_s=S_sH$ using
Lemma~\ref{lem:finite-order-smoothing}.  Then, for every fixed $m$,
\begin{align}
 |j^2G_s|_{m,s}+|H_s|_{m,s}&\le C_{m,r}B,
 \label{eq:exact-smoothed-bounds}\\
 |j^2(G-G_s)|_{m,\tau}+|H-H_s|_{m,\tau}
 &\le C_{m,r}\left[
 B(s/t)^r+\mathcal S_m(t;G,H)(\tau/t)^r\right].
 \label{eq:exact-smoothing-tail}
\end{align}
For the first inequality use \eqref{eq:smooth-derivative} with $q=2$
and $q=0$, respectively.  To prove the second, split the derivatives
at $r$.  A derivative of order $j\le r$ of $j^2G$ or $H$ contributes
at most
\[
 C B t^{-r}s^{r-j}\tau^j
 =C B(s/t)^r(\tau/s)^j
 \le C B(s/t)^r.
\]
For $j>r$, boundedness of smoothing in the original derivative norm
gives a contribution at most
\[
 C_m(\tau/t)^j\mathcal S_m(t;G,H)
 \le C_m(\tau/t)^r\mathcal S_m(t;G,H).
\]
In particular, the constants in \eqref{eq:exact-smoothed-bounds} do
not use any derivative of the true input above order $r+2$; the
possibly large higher derivatives occur only linearly in
\eqref{eq:exact-smoothing-tail}.

\subsection{A recurrence at every derivative order}

Fix a small $z$ for the moment.  Lemma~\ref{lem:polynomial-phase-data}
provides phase data for $f_z$ with reference metric $h_*$.  They remain
admissible in the $C^2$ ball of radius
\[
 \rho_z=z^4
\]
about $f_z$.  Their bounds in every fixed order are polynomial in
$z^{-1}$.  Put $B=z^{-4}$.  The current integer $k\ge10$ controls the
amplitude and the number of derivatives to be retained.  We use
\begin{equation}\label{eq:exact-scales}
 \beta=\frac65,\qquad \lambda=\frac{11}{10},\qquad
 \tau=t^\beta,\qquad s=t^\lambda,\qquad
 \delta=t^k,\qquad r(k)=q(k)=40(k+1).
\end{equation}
Here $q(k)$ is the accuracy order in
Proposition~\ref{prop:small-increment}.  The integer $r(k)$ is the
order controlled at the current step; the estimates below will also hold
at every fixed observation order $m>r(k)$, without assuming a bound there.

For a current map $G$, define the normalized defect
\begin{equation}\label{eq:exact-normalized-defect}
 H=\delta^{-2}(h_*-\gamma(G)),\qquad
 \mathcal S_m=\mathcal S_m(t;G,H).
\end{equation}
The step will be applied under the finite-order hypotheses
\begin{equation}\label{eq:exact-step-hypotheses}
 \|G-f_z\|_{C^2}\le\rho_z/2,\qquad
 \|H-h_*\|_{C^0}<\varepsilon_*,\qquad
 \mathcal S_{r(k)}\le B,
\end{equation}
where $\varepsilon_*>0$ is a fixed sufficiently small radius for the
positive tensor decompositions.

\begin{lemma}[One smoothing step]\label{lem:exact-step}
Under \eqref{eq:exact-step-hypotheses}, take either
$\delta'=\tau^k$ or $\delta'=\tau^{k+1}$.  For sufficiently small
$t$, there is a smooth increment $U$ such that, with
\[
 G'=G+U,\qquad H'=\delta'^{-2}(h_*-\gamma(G')),
 \qquad \mathcal S'_m=\mathcal S_m(\tau;G',H'),
\]
for every fixed $m\ge0$,
\begin{align}
 |U|_{m,\tau}&\le C_m(z,k)\delta\tau,
 \label{eq:exact-increment-bound}\\
 \mathcal S'_m
 &\le \|G\|_{C^2}+\|h_*\|_{C^0}
       +C_m(z,k)t^{1/5}(1+\mathcal S_m).
 \label{eq:exact-all-order-recurrence}
\end{align}
The scale threshold and the constants are independent of derivatives
of the true inputs $G,H$ above the controlled orders $r(k)+2,r(k)$.
For fixed $k$, the inverse threshold is polynomially bounded in $z^{-1}$;
for fixed $k,m$, the constants have the same property.
\end{lemma}

The same construction therefore has two different uses.  At the controlled
order $r(k)$, a small scale preserves its bound.  At any higher fixed
order, the current norm may be large, but it is multiplied by a coefficient
that tends to zero.  The latter fact will let us increase $k$ after a
finite number of steps.

\begin{proof}
Smooth with moment order $r(k)$.  By
\eqref{eq:exact-smoothing-tail}, $G_s$ belongs to the admissible
$C^2$ ball for small $t$.  The tensor
\[
 H_s-(\delta'/\delta)^2h_*
\]
belongs to the required neighborhood of $h_*$: the smoothing tail
tends to zero, and $\delta'/\delta\to0$.  Apply
Proposition~\ref{prop:small-increment} in its unperturbed case to this
tensor at $G_s$, with amplitude $\delta$ and accuracy $q(k)$.  Write
the resulting increment as $U$.  The proposition gives
\eqref{eq:exact-increment-bound} and
\begin{equation}\label{eq:exact-mode-error}
 |E|_{m,\tau}\le C_m(z,k)
 \left[\delta(\tau/s)^{q(k)}+\delta^3/\tau\right],
\end{equation}
where
\[
 E=\gamma(G_s+U)-\gamma(G_s)
   -\delta^2\bigl[H_s-(\delta'/\delta)^2h_*\bigr].
\]
The constants have the stated independence from higher true-input
norms because \eqref{eq:exact-smoothed-bounds} bounds every required
weighted derivative of $G_s,H_s$.  For fixed orders they are bounded
by fixed powers of $z^{-1}$, by the polynomial phase bounds and the
polynomial dependence in Remark~\ref{rem:small-polynomial-control}.

The increment was constructed using $G_s,H_s$, but we add it to the
unsmoothed map $G$.  To compare the two metric increments, define the
symmetric tensor
\[
 \mathcal C(V,W)_{ij}
 =\langle\partial_iV,\partial_jW\rangle
  +\langle\partial_jV,\partial_iW\rangle.
\]
Since the induced metric is quadratic in the first derivatives,
\begin{equation}\label{eq:exact-error-identity}
 \delta'^2(H'-h_*)
 =\delta^2(H-H_s)-E-\mathcal C(G-G_s,U).
\end{equation}
In particular, there is no square of the smoothing tail.  From
\eqref{eq:exact-increment-bound}, $|dU|_{m,\tau}\le C_m(z,k)\delta$;
\eqref{eq:exact-smoothing-tail} bounds the other factor in the last
term.  Hence the high derivatives of the true input enter the new
defect only linearly.

For completeness, the powers needed to divide
\eqref{eq:exact-error-identity} by $\delta'^2$ are displayed below.
The smaller choice $\delta'=\tau^{k+1}$ is the more demanding one.
Using \eqref{eq:exact-scales}, the respective exponents of $t$ are
\begin{equation}\label{eq:exact-exponent-table}
\begin{array}{c|c}
 \text{term}&\text{exponent after division by }\delta'^2\\ \hline
 \delta(\tau/s)^{q(k)}&(13k+8)/5\\
 \delta^3/\tau&(3k-18)/5\\
 \delta^2 B(s/t)^{r(k)}&(18k+8)/5\\
 \delta^2\mathcal S_m(\tau/t)^{r(k)}&(38k+28)/5\\
 \delta B(s/t)^{r(k)}&(13k+8)/5\\
 \delta\mathcal S_m(\tau/t)^{r(k)}&(33k+28)/5.
\end{array}
\end{equation}
The factors $B$ are included in the constants for fixed $z$.
Every exponent is positive for $k\ge10$.  We obtain, in particular,
\begin{equation}\label{eq:exact-defect-recurrence}
 |H'-h_*|_{m,\tau}
 \le C_m(z,k)t^{1/5}(1+\mathcal S_m)
 \qquad(m\ge0).
\end{equation}
This weaker common exponent will also control the map derivatives.
Indeed, inherited derivatives of $G$ of positive weighted order gain
at least $\tau/t=t^{1/5}$, while
\[
 |j^2U|_{m,\tau}
 \le C_m(z,k)\delta/\tau
 =C_m(z,k)t^{k-6/5}.
\]
The fixed tensor satisfies
$|h_*|_{m,\tau}\le\|h_*\|_{C^0}+C_m\tau$.
Together these estimates give
\eqref{eq:exact-all-order-recurrence}.  Its constants use only the
smoothed-input bounds, while all higher true-input derivatives have
remained in the displayed linear term $\mathcal S_m$.
The polynomial bounds already noted for the phase data and finite mode
formulas, with the strict positive powers of $t$ in
\eqref{eq:exact-exponent-table}, give the stated polynomial thresholds.
\end{proof}

It remains to initialize this step from the finite assumptions of
Proposition~\ref{prop:exact-correction}, keep the maps in the allowed
$C^2$ ball, and schedule increases of $k$.

\subsection{Initialization and passage through the derivative orders}

Because $\|f_z\|_{C^2}=O(z^{-2})$, for small $z$ the distance condition
in \eqref{eq:exact-step-hypotheses} implies
\begin{equation}\label{eq:exact-baseline}
 \|G\|_{C^2}+\|h_*\|_{C^0}<B/4.
\end{equation}
For each fixed $z,k$, choose a positive scale threshold that makes
Lemma~\ref{lem:exact-step} applicable, preserves the defect neighborhood,
and, by \eqref{eq:exact-all-order-recurrence} with $m=r(k)$, preserves
$\mathcal S_{r(k)}\le B$.  Impose also
\begin{equation}\label{eq:exact-summable-increment}
 \|U\|_{C^{\lfloor k/2\rfloor}}\le\rho_z t.
\end{equation}
This last threshold is available because
\eqref{eq:exact-increment-bound} gives the upper bound
\[
 C(z,k)t^{k+(6/5)(1-\lfloor k/2\rfloor)},
\]
and its exponent is greater than one for every $k\ge10$.
The thresholds can cover both choices of $\delta'$ at once.

For the initial block $k=10$, only finitely many orders and
inequalities have been used.  Their constants are polynomial in
$z^{-1}$, and every smallness condition has a strict positive power
of $t$.  Therefore there is a fixed $\mu>1$ such that all these
conditions hold for $t\le z^\mu$ and sufficiently small $z$.
Enlarge $\mu$ if necessary, and set
\[
 R=r(10)=440,\qquad N>20\mu+1,
\]
with $N$ an integer.  These choices use polynomial exponents, not the
implicit constants in the assumptions of the proposition.

Start at $t=z^\mu$, $k=10$, $\delta=t^{10}$, and
\[
 G=\sqrt{1-\delta^2}\,f_z.
\]
Its normalized defect has the exact form
\[
 H=h_*+(1-\delta^{-2})(\gamma(f_z)-h_*).
\]
It therefore tends to $h_*$ in $C^R$.  The weighted derivatives of
$G$ through the required order are bounded by
\[
 C z^{-2}\sum_{j=0}^R z^{(\mu-1)j},
\]
so $\mathcal S_R<B/2$ for small $z$.  The initial displacement is
$O(z^{20\mu-2})$, which is at most $\rho_z/4$ for sufficiently small
$z$, since $\mu>1$ and $20\mu-2>4$.  This verifies all the initial conditions.

At every subsequent step replace $t$ by $t^{6/5}$.  Thus, if $t_n$
is the scale at step $n$,
\[
 t_n=t_0^{(6/5)^n},\qquad \sum_{n=0}^\infty t_n<1/4
\]
for small enough $t_0$.  Equations
\eqref{eq:exact-summable-increment} and the initial displacement keep
all maps within $\rho_z/2$ of $f_z$ in $C^2$.  Hence the phase data
and \eqref{eq:exact-baseline} remain valid.

We explain why no derivative order prevents the construction from
advancing.  Suppose that $k$ is held fixed and let $m$ be any fixed
integer.  Apply Lemma~\ref{lem:exact-step} repeatedly with
$\delta'=\tau^k$.  Its all-order recurrence says that the numbers
$x_n=\mathcal S_m(t_n;G_n,H_n)$ satisfy
\[
 x_{n+1}\le B/4+c_n(1+x_n),\qquad
 c_n=C_m(z,k)t_n^{1/5}\longrightarrow0.
\]
Every $x_n$ is finite because all finite-stage maps are smooth.  Once
$c_n\le1/4$, this recurrence first bounds the sequence; subsequently
$c_n\to0$ gives $\limsup x_n\le B/4$.  In particular, after finitely
many further steps $x_n\le B$.

Apply this observation with $m=r(k+1)$.  Wait in the block with
integer $k$ until that higher bound holds and the scale is below the
threshold for the block with integer $k+1$.  Wait also until
\[
 C_{r(k+1)}(z,k)t^{1/5}(1+B)\le B/2,
\]
where the constant is the one for the current block in
\eqref{eq:exact-all-order-recurrence}.  All three conditions hold after
finitely many steps.  Now choose $\delta'=\tau^{k+1}$.  That recurrence,
already proved for this smaller target amplitude, bounds the new higher
norm by $B/4+B/2<B$.  The new scale $\tau<t$ also satisfies the next
block's threshold.  Thus each block has finite length and $k$ tends to
infinity.

For every fixed $m$, all sufficiently late blocks satisfy
$\lfloor k/2\rfloor\ge m$.  Equation
\eqref{eq:exact-summable-increment} then makes the tail of the
increments absolutely summable in $C^m$.  The maps converge in every
$C^m$ to a smooth map $\widetilde f_z$.  Their defects tend to zero in
$C^0$, since $\delta=t^k\to0$ and $H$ stays in a fixed $C^0$
neighborhood of $h_*$.  Continuity of the induced metric in $C^1$
yields $\gamma(\widetilde f_z)=h_*$.  Finally,
\[
 \|\widetilde f_z-f_z\|_{C^2}
 \le\rho_z/4+\rho_z\sum_{n=0}^\infty t_n
 \le\rho_z/2\longrightarrow0.
\]
This proves Proposition~\ref{prop:exact-correction}.

\begin{remark}
The finite assumptions control the first block uniformly by powers
of $z$.  Derivatives above order $R+2$ may be arbitrarily large; at a
fixed $z$ they affect only the finite waiting times in later blocks.
The construction gives $C^2$ proximity on the whole surface.  It does
not assert that the exact correction vanishes outside a prescribed
support.
\end{remark}
